
\input texsis
\book
\overfullrule=0pt
\def\frac#1#2{#1\over #2}
%
\def\bmatrix#1{\left[ \matrix{#1} \right]}
\def\be{{\beta}}
\def\toone{{t+1}}

%\def\frac#1#2{#1\over #2}
%\magnification=1200
\overfullrule=0pt

\centerline{\bf RMT5 exercises}
\medskip



\noindent{\bf 1.} \quad {\it Harrison-Kreps (1978) 101}
\medskip

\noindent One share of a risk-free asset pays a constant dividend $d >0$ at each time $t = 0, 1, 2, \ldots$.
The asset is traded by two types of risk neutral investors that we call  `temporarily odd' and `temporarily even'. There are equal numbers of the two types.
Both types are infinitely lived.  There is one share of the asset for each pair of agents.  Each type of agent has ``deep pockets'' and can afford to hold the total supply of the asset.
A temporarily odd investor
has a sequence of one-period discount factors $\beta_o, \beta_e, \beta_o, \beta_e, \ldots $, while a temporarily  even investor
has a sequence of one-period discount factors   $\beta_e, \beta_o, \beta_e,  \beta_o, \ldots $, where $\beta_e > \beta_o$
and $\beta_i \in (0,1)$ for $i=o, e$. For example, if a temporarily odd investor were to hold the asset forever, he would value
it at $p_o = d( 1 + \beta_o + \beta_o \beta_e + \beta_o \beta_e \beta_o + \cdots )$.

\medskip
\noindent {\it Personal  valuation of dividend stream}

\medskip
\noindent   Note that if he were to hold it forever, a temporarily odd investor would value the asset at
$$ p_o = d + \beta_o p_e$$
where
$$ p_e = d + \beta_e p_o .$$
If he were to hold it forever, a temporarily even investor would value the asset at
$p_e$.

\medskip
\noindent {\bf a.} Verify that $p_o = {\frac{1 + \beta_o}{1-\beta_e \beta_o}} d $ and
$ p_e = {\frac{1 + \beta_e}{1-\beta_e \beta_o}} d $.

\medskip
\noindent {\bf b.} Verify that $p_e > p_o$.


\medskip
\noindent{\bf c.} Please imagine a  ``perpetually even agent'' who has a constant one-period discount factor equal to
$\beta_e$.  (There is actually no perpetually even agent in our model, so this is just a mental exercise.)  Please verify
that a perpetually even agent would pay $\bar p_e = {\frac{1}{1-\beta_e}}d $ for the asset.



\medskip
\noindent {\bf d.}  Now assume that the  two actual types of agents, each of which oscillates between being temporarily odd and temporarily even,  trade the asset
under the following rules. They cannot short sell the asset, so their holdings of the asset must be nonnegative.  Ownership of the asset at the beginning of a period entitles the owner to sell the asset ``cum dividend'', meaning that the buyer
receives the current dividend and the right to sell the asset at the beginning of the next period.
Argue  that the agent who owns  the asset each period after trading  is the one who is temporarily even in that period.  Write a Bellman equation for the equilibrium price
of the asset under that assumption.
At what price does the asset trade? What is trading volume each period?  Please describe who trades what when.

\medskip
\noindent{\bf e.}  How does the equilibrium price of the asset compare with the amount that would be paid by a perpetually even agent?

\medskip
\noindent{\bf f.} How does the equilibrium price of the asset compare with the personal value  assigned  to the stream of dividends  by each type of agent?
{\bf Hint:} Please verify  that ${\frac{\bar p_e}{p_e}} = {\frac{ 1- \beta_o \beta_e}{1 - \beta_e^2}} >1$ and tell how this calculation bears on the answer.

\medskip
\noindent{\bf g.}  Please explain whether you agree or disagree with the following statement.  ``The asset trades for more than anyone personally thinks that the asset's dividend stream is worth because
acquiring the asset now gives a purchaser this period  the option to sell it next period to someone who is more patient than  he or she will be.''







\bye






\bye
